% !PRINTED 359
\weberpage{359}
If the double tangents fall together into one line, two double points appear.

On the other hand, without singular points arising, the two points of contact of a double tangent may very well fall together, and the double tangent thus pass over into a tangent with four-point contact. This occurs when the line $x_1$ touches the conic $u$.

The form (2), too, can still be produced in infinitely many ways while retaining $x_1, y_1$.

To see this, let us assume that
\[
f = x_1 y_1 v - u^2 = x_1 y_1 v_1 - u_1^2,
\]
from which the identical equation follows:
\weq{3}{x_1 y_1 (v - v_1) = (u - u_1)(u + u_1).}
Now if $u - u_1$ were divisible by $x_1$, and $u + u_1$ by $y_1$, then at the point of intersection of $x_1$ and $y_1$, $u$, and consequently $f$, would also vanish. This point of intersection would lie on the curve $f$, which, so long as the curve has no singular point, is not possible. Hence one of the two factors $u - u_1$, $u + u_1$ in (3) must be divisible by $x_1 y_1$. Since the sign of $u_1$ is not yet determined, we may assume that $u - u_1$ is divisible by $x_1 y_1$, and therefore identical with $x_1 y_1$ up to a constant factor. Then, if $\lambda$ is such a factor, we have
\[
u_1 = u + \lambda x_1 y_1,
\]
and the function $f$ takes the form:
\weq{4}{f = x_1 y_1 (v + 2\lambda u + \lambda^2 x_1 y_1) - (u + \lambda x_1 y_1)^2.}
Conversely, for every arbitrary $\lambda$ the formulae (2) and (4) are identical with one another.

Now if we can determine $\lambda$ so that the quadratic function $v + 2\lambda u + \lambda^2 x_1 y_1$ splits into two linear factors, then, setting $u + \lambda x_1 y_1 = u_1$, $f$ takes by (4) the shape
\weq{5}{f = x_1 y_1 x_2 y_2 - u_1^2,}
and this form shows that $x_2, y_2$ are \gesperrt{two further double tangents}, and that the eight points of contact of the double tangents $x_1, y_1, x_2, y_2$ lie on a conic $u_1 = 0$.
